\documentclass[aspectratio=169]{beamer}
\usepackage[utf8]{inputenc}
\usepackage{amsmath, amssymb}
\usepackage{graphicx}
\usepackage{tcolorbox}
\definecolor{UNIBBlue}{RGB}{0, 51, 102}
\definecolor{UNIBYellow}{RGB}{255, 204, 0}
\setbeamercolor{palette primary}{bg=UNIBBlue,fg=white}
\setbeamercolor{title}{fg=UNIBBlue}
\setbeamercolor{frametitle}{bg=UNIBBlue,fg=white}
\title{Optimisasi untuk Teknik Elektro}
\subtitle{Minggu 6: Konsep Konveksitas}
\author{Ir. Novalio Daratha S.T., M.Sc., Ph.D.}
\date{Semester Ganjil 2026}
\begin{document}
\begin{frame}\titlepage\end{frame}
\begin{frame}{Tujuan Pembelajaran}
\begin{itemize}
    \item Mendefinisikan himpunan cembung dan fungsi cembung.
    \item Mengidentifikasi apakah suatu masalah optimisasi bersifat cembung.
    \item Menjelaskan mengapa optimisasi cembung sangat penting dalam keteknikan.
\end{itemize}
\end{frame}
\begin{frame}{Himpunan dan Fungsi Cembung}
\textbf{Himpunan Cembung:} Garis yang menghubungkan dua titik mana pun dalam himpunan tersebut sepenuhnya berada di dalam himpunan.\\
\textbf{Fungsi Cembung:} Garis (secant) yang menghubungkan dua titik pada kurva selalu berada di atas atau pada kurva fungsi tersebut.
\end{frame}
\begin{frame}{Pentingnya Konveksitas}
\begin{tcolorbox}[colback=yellow!5,title=Sifat Emas Konveksitas]
Setiap minimum lokal pada fungsi cembung dalam himpunan cembung adalah \textbf{minimum global}.
\end{tcolorbox}
Dalam teknik elektro, mendesain sistem komunikasi atau filter sinyal seringkali dapat diformulasikan sebagai masalah konveks (seperti Convex Optimization, SOCP). Jika model kita konveks, kita dijamin menemukan solusi optimal global secara efisien.
\end{frame}
\end{document}